% !PRINTED 88
\weberpage{88}
%<<CONT
Wenn der imaginäre Teil von $\omega$ ins Unendliche wächst, so
daß der absolute Wert von $q$ verschwindet, so wird
\[
\vartheta_{00}(v) = 1, \quad \vartheta_{01}(v) = 1, \quad q^{-\frac{1}{4}}\vartheta_{10}(v) = 2\cos\pi v,
\]
\[
q^{-\frac{1}{4}}\vartheta_{11}(v) = 2\sin\pi v.
\]
Nehmen wir $\omega$ rein imaginär, also $q$ reell, positiv und echt
gebrochen an, so sind für ein reelles $v$:

1. $\vartheta_{00}(v)$, $\vartheta_{01}(v)$, $\vartheta_{10}(v)$, $\vartheta_{11}(v)$ reell, und wenn $v$ zwischen 0
und $\frac{1}{2}$ liegt, positiv [nach § 24 (7)], folglich auch $\vartheta_{00}$, $\vartheta_{01}$, $\vartheta_{10}$, $\vartheta'_{11}$
positiv;

2. $\vartheta_{00}(iv)$, $\vartheta_{01}(iv)$, $\vartheta_{10}(iv)$, $-i\vartheta_{11}(iv)$ reell, und solange $v$
zwischen 0 und $\dfrac{-i\omega}{2}$ liegt, positiv; ferner mit Zuziehung der
Formeln § 21, (8)

3. $\vartheta_{00}(\frac{1}{2}+iv)$, $\vartheta_{01}(\frac{1}{2}+iv)$, $i\vartheta_{10}(\frac{1}{2}+iv)$, $\vartheta_{11}(\frac{1}{2}+iv)$
reell, und solange $v$ zwischen 0 und $\dfrac{-i\omega}{2}$ liegt, positiv;
\[
4. \quad q^{\frac{1}{4}}e^{\pi iv}\vartheta_{00}\left(\frac{\omega}{2}+v\right), \quad -iq^{\frac{1}{4}}e^{\pi iv}\vartheta_{01}\left(\frac{\omega}{2}+v\right),
\]
\[
q^{\frac{1}{4}}e^{\pi iv}\vartheta_{10}\left(\frac{\omega}{2}+v\right), \quad -iq^{\frac{1}{4}}e^{\pi iv}\vartheta_{11}\left(\frac{\omega}{2}+v\right)
\]
reell, und solange $v$ zwischen 0 und $\frac{1}{2}$ liegt, positiv.

\begin{center}
\textbf{§ 26. \gesperrt{Entwickelung von $\vartheta$-Quotienten.}}
\end{center}

Bedeutet $v$ eine Variable und $a$ eine beliebige Konstante, so
sind die Quotienten
\[
\frac{\vartheta_{g_1,g_2}(v+a)}{\vartheta_{g'_1,g'_2}(v)}
\]
doppeltperiodische Funktionen zweiter Art und erster Ordnung
mit den Periodizitätsfaktoren $(-1)^{g_1+g'_1}$, $(-1)^{g_2+g'_2}e^{-2\pi ia}$, und
wenn man einen Exponentialfaktor $e^{\lambda v}$ hinzufügt, so kann man $\lambda$
und $a$ so bestimmen, daß die Periodizitätsfaktoren beliebig gegebene Größen werden. Indem man die Hauptcharakteristiken
$(g_1,g_2), (g'_1,g'_2)$ auf alle mögliche Arten wählt, erhält man 16 solcher
Funktionen. Wir gehen aus von einer unter ihnen, für die wir
unter Hinzufügung eines konstanten Faktors
\weq{1}{F = \frac{\vartheta'_{11}\vartheta_{11}(v+a)}{2i\pi\vartheta_{11}(v)\vartheta_{11}(a)}}
%CONT>>
